\begin{annexitem}{annex:p12.n-c-035.5.h-minimal-instance}{thm:p12-eu-h-minimal-instance}{Both left-occurring and right-only generating images}
\begin{proof}
Order applicable ground pairs lexicographically. A nonempty set has
a least pair because the ground term order is total and well founded.
All rule presentations are restricted to their identity variables.
Nontriviality rules out an oriented variable-left instance: if the
left variable is absent on the other side, the identity makes every
model a singleton; if present in a proper other term, its ground
image occurs properly there and cannot be greater. Otherwise the
identity is trivial. A proper ground context is greater than its
subterm, since the contrary direction, iterated by monotonicity,
would produce an infinite descending context chain.

If a generating image \(\rho(v)\) is reducible and v occurs in
the nonvariable left pattern, its redex has a left side smaller
than \(l\rho\), because it is a proper subterm. It is also
applicable in the ambient term. This contradicts minimality of
the pair's first coordinate. If v occurs only on the right,
perform one reduction in its image, producing \(\rho'\).
The left side is unchanged; replacing all right occurrences
strictly decreases \(r\rho\) by monotonicity, while the new
identity instance remains oriented. This contradicts minimality
of the second coordinate. Thus every image is R-normal, and
the pair belongs to \(R_E\).

An R-normal term is \(R_E\)-normal since \(R_E\subseteq R\).
Conversely an R-reducible finite term has such a minimal applicable
pair and therefore has an \(R_E\)-redex. This proves equality
of the normal terms, also for a nonground term whose redex is a
ground subtree. No procedure searching an infinite rule set is
asserted; the least-pair choice is the proof's well-founded choice.
\end{proof}
\end{annexitem}

\begin{annexitem}{annex:p12.n-c-035.5.h-basic-normalization}{thm:p12-eu-h-basic-normalization}{The source basic skeleton is constructed at every step}
\begin{proof}
R terminates by the ordering. Start with skeleton s and substitution
\(\theta\). While the instantiated term is reducible, choose a
minimal applicable pair as in the preceding lemma. Its generating
substitution \(\rho\) is reduced, and its identity is chosen
fresh for the current skeleton and the finite substitution data.
The redex cannot be at or below a variable-image occurrence:
those images are \(R_E\)-normal by hypothesis. Hence it occurs
at a nonvariable position in the skeleton. Replace that occurrence
by the fresh right pattern r, and extend the accumulated map to
\(\theta\rho\). Freshness makes old images unchanged; every
new image is normal. The instantiated successor is exactly the
chosen ground rewrite, by the finite-term composition law.
Repeat this argument with the actual updated skeleton, not the
original one. All steps are based, and all used instances are in
\(R_E\). Termination makes this sequence finite.

Its terminal term is \(R_E\)-normal and thus R-normal by the
preceding normal-term equivalence. Ground confluence makes it
the same endpoint as every R-normalization of the initial term.
For the two-sided assertion, E-equal ground endpoints have the
same R-normal form: ground every intermediate variable in a finite
equational derivation by a fixed ground term, fixing the endpoints,
then orient its ground identity steps using the total order.
Confluence and termination give their common normal form. Apply
the constructed basic normalizations to both sides with fresh
rule variants. Conversely such a basic proof is an equational
conversion, since every used instance is an E-identity.
R-reduced substitutions are \(R_E\)-reduced, so the stated
corollary follows without assuming the basic proof in advance.
\end{proof}
\end{annexitem}

\begin{annexitem}{annex:p12.n-c-035.5.h-basic-transport}{thm:p12-eu-h-basic-transport}{Finite substitution transport of the actual position sequence}
\begin{proof}
The first instantiated term is unchanged:
\((q\mu)\alpha=q(\mu\alpha)=q\alpha\).
An eligible nonvariable position of q retains all its constructor
ancestors under \(\mu\); substitution only adds nodes below
variable leaves. Thus the old redex remains at a nonvariable
position of \(q\mu\). The fresh right pattern of this step
contains no \(\mu\)-domain variable, so replacing at that
position commutes with \(\mu\). Extend the ground witness
by its same fresh reduced generating map. Ground old-variable
images and freshness preserve \(\mu\alpha=\alpha\) on
every relevant older variable. The next ground term is again
identical to the old proof's next term. Induct through the finite
position sequence, renaming future variants away from the finite
\(\mu\) data. Every transported step uses the same ground
\(R_E\) instance and remains based on its updated skeleton.
The endpoint and length are unchanged. Newly expanded variable
regions have ground value equal to the old variable's normal
image, so this argument has not used an unproved reduction inside
an inserted image. Minimum length in the enlarged skeleton
representation is at most the length of this transported proof.
\end{proof}
\end{annexitem}

\begin{annexitem}{annex:p12.n-c-035.5.h-minimum-blocks}{thm:p12-eu-h-minimum-blocks}{The source tuple with reduced minima and full finite matching blocks}
\begin{proof}
Each equation has a finite basic \(R_E\) proof by the preceding
theorem. Its proof lengths are a nonempty subset of natural
numbers and have a minimum. Choose a minimum independently for
each equation occurrence. Keep the reduced ground witness,
literal solution of S, and the reachable-state invariant that
\(\operatorname{Dom}(\sigma_S)\) is absent from P. Fresh LP
variants and the actual removing-binding U transitions preserve
this last invariant; it is not assumed for arbitrary unrelated
P,S pairs. The initial S is empty.

For an equation with equal normal instantiated sides, one actual
nonfailure U step is enabled. Different constructors or a proper
occurs obstruction are impossible for its literal finite ground
witness. Removing a variable binding retains that same witness
because \(\delta\theta=\theta\), and updates old solved
images while preserving solved form. Other equation instances
are unchanged. Their former basic proofs transport by the
preceding lemma, so \(m_{R_E}\) cannot increase. New children
of a decomposed normal equation are equal normal terms and
contribute zero. If the first coordinate strictly falls, descent
is established. Otherwise the source three counters decrease
exactly as in the syntactic helper: elimination removes its
variable from all P; deletion/decomposition lower symbol count
if variable count is unchanged; orientation lowers the last
count if both earlier counts stay fixed. The normal-equation
classification of other constraints cannot change because their
ground instances have not changed.

If no such equation is selected, take the first step on a
nonempty leg of one chosen minimum \(R_E\) proof. It uses
\(l\rho\to r\rho\), with reduced generating \(\rho\).
The based position is a nonvariable skeleton position and l
has the same root. Extend the witness by this fresh \(\rho\),
leaving all old values unchanged. LP produces the replacement
equation and the matching equation \(l=^{?}u\), literally
solved by that extended witness. Perform its mandatory first
Decomposition, then solve the remaining matching constraints
completely with the actual U helper. This is finite and cannot
fail, since the literal witness exists. If identical constructors
would have been removed by Trivial, Decomposition followed by
child Trivial steps gives the same output while respecting LP.
The helper's global substitution \(\mu\) satisfies
\(\mu\alpha=\alpha\) for the extended witness \(\alpha\).
All new bindings have normal instantiated images, because their
left variables have those reduced witness images. Thus the
extended witness remains reduced and literally solves the updated S.

The matching equations have disappeared and contribute no proof
length. Remove the first step from the selected old basic proof;
its remaining legs have total length one less. Transport them
through \(\mu\) by the preceding position-by-position lemma.
Transport every other equation's old minimum proof in the same
way. They supply basic \(R_E\) proofs of the successor
whose total length is at most the old \(m_{R_E}-1\).
New minima can only be shorter. Hence the full LP/matching block
strictly decreases the first coordinate, regardless of growth
of the syntactic counters or addition of fresh variables.

The lexicographic tuple is well founded. Induction therefore
constructs a finite sequence of these actual steps/blocks ending
with P empty. Its S is solved, and the retained witness gives
\(\sigma_S\alpha=\alpha\) on the fixed input domain.
Fresh extensions never changed those values, so restriction
gives \(x\theta=x\sigma_S\lambda\) on the original X.
The compared domain is not silently changed to the shrinking
variable set of a later P. For nonground original parameters,
use the already declared fresh-constant domain and unfreeze
the finite witness at the end.

The three examples establish why the printed raw inference
cannot be substituted for this proof. A matching identity at
a reducible term needs two normalization legs, not a reflexive
normal-form proof. A single decomposition can leave a reducible
matching constructor pair. In the erasing example, the one-step
basic unrestricted proof has a nonnormal generating image;
the reduced minimum is instead two. Thus both corrections do
real mathematical work and do not merely rename the desired
completeness statement.
\end{proof}
\end{annexitem}

\begin{annexitem}{annex:p12.n-c-035.5.h-convergent-hereditary}{thm:p12-eu-h-convergent-hereditary}{Full overlap agreement and actual substitution-history marks}
\begin{proof}
At each LP in the just constructed trace, \(l\rho=u\theta\)
literally. If l and u both have constructor heads at a common
position, simultaneous substitution cannot change those heads.
Literal equality therefore forces the same symbol and arity
there. This verifies the whole nonvariable overlap test, not
only the original root test, at every actual LP choice.

Give every constructor occurrence introduced by Variable
Elimination in P an ineligible history mark, and inherit marks
through decomposition and later copying. Such a substituted
term t has ground value \(t\alpha=x\alpha\), because the
selected binding is literally solved by the witness. The latter
value is normal. All old-variable ground values remain unchanged
in subsequent witness extensions, so the marked region stays
normal for the remainder of the trace. Every newly marked copy
has the same justification. Syntactic constraint solving may
decompose or remove a marked term; only LP is restricted.

An \(R_E\)-redex is also an R-redex. It cannot occur anywhere
inside one of these normal marked regions, including at its
root. The first step of every selected minimum basic proof
therefore occurs at an eligible historical position. New
identity-template occurrences introduced by that LP are not
substitution copies; their later normality/mark status is
determined by the same literal-solve calculation. Hence all
LP choices satisfy both restrictions, while all finite U
solves remain legal. The previous tuple induction and final
factorization are unchanged. This proves the stated convergent
regime, not the unproved transport of these marks through an
arbitrary sequence of nested critical-pair simulations.
\end{proof}
\end{annexitem}

